\subsection{Numerical implementation}
\label{sec:numerics}

The model is solved with a finite-volume code built on the block-structured adaptive mesh refinement framework AMReX \cite{amrex2019}. The domain \(0\le x\le L\) is divided into cells, and the cell averages of \(n_e\), \(n_i\), \(n_n\), \(\mathcal{E}\), \(T_h\), \(\rho u\) and \(\phi\) are the discrete unknowns. Every conservation law is written in flux form, with the fluxes evaluated on the cell faces. The species fluxes~(\ref{eq:flux-e}) and (\ref{eq:flux-i}) are split into a density-diffusion part, \(-D\,\partial n/\partial x\), and a non-diffusive part, \(nv\), where the velocity \(v\) collects the drift in the field, the temperature-gradient and thermal-force terms, the mixture force and the collisional drag. In the sheath, the drift and diffusion fluxes of the electrons nearly cancel, and their sum is two to three orders of magnitude smaller than either. A small relative error in either part would therefore be a large error in the net flux. The face flux is accordingly the exponentially fitted (Scharfetter--Gummel) flux \cite{scharfetter1969},
\begin{equation}
\Gamma_{f}=\frac{D_f}{\Delta x}\Big[B(-\psi)\,n_{\rm L}-B(\psi)\,n_{\rm R}\Big],
\qquad
\psi=\frac{v_f\,\Delta x}{D_f},
\qquad
B(\psi)=\frac{\psi}{e^{\psi}-1},
\label{eq:sg}
\end{equation}
where L and R denote the cells on either side of the face and \(\psi\) is the cell Péclet number. Equation~(\ref{eq:sg}) is exact for the exponential profile of a Boltzmann balance; it reduces to central differencing as \(\psi\to0\) and to upwinding as \(|\psi|\to\infty\). The enthalpy part of the electron energy flux~(\ref{eq:heatflux-e}) is computed from this discrete electron flux, \(\tfrac52k_BT_e\,\Gamma_{e,f}\), so that each electron crossing a face carries exactly its own enthalpy. Heat conduction in both energy equations and Poisson's equation~(\ref{eq:poisson}) are discretized with central differences, and the face field \(E_f=-(\phi_{\rm R}-\phi_{\rm L})/\Delta x\) is used in every face flux and in the Joule heating. The heavy-gas momentum equation~(\ref{eq:heavy-momentum}) is discretized with a Roe-type flux \cite{roe1981}. At the electrode, the wall fluxes~(\ref{eq:bc-electron})--(\ref{eq:bc-neutral}) are imposed as the face fluxes, evaluated with the densities and temperatures of the cell adjacent to the wall; at \(x=L\), the Dirichlet values of Section~\ref{sec:model-bc} are imposed on the boundary face.

All cases in this study use a uniform base mesh of 512 cells with four levels of refinement towards the electrode, each with a refinement ratio of 4. The refined regions are nested layers adjacent to the electrode, \(1.0\times10^{-4}\), \(2.0\times10^{-6}\), \(7.3\times10^{-7}\) and \(2.7\times10^{-7}\)~m deep, with cell sizes of 244, 61, 15.3 and 3.8~nm, so that the mesh has 944 cells in total. These sizes are chosen to resolve the length scales estimated in Section~\ref{sec:scales involved}, down to a fraction of the Debye length, at both the cathode and the anode. The levels are coupled conservatively. The equations are advanced in time, from a uniform initial state at the bulk-plasma conditions or from a neighboring steady state, until a steady state is reached. All equations, including Poisson's equation and the galvanostatic condition~(\ref{eq:bc-current}), are solved together at each time step. A run is considered steady when, for each of \(\phi\), \(n_e\), \(n_i\), \(\mathcal{E}\) and \(T_h\), the relative change over a time step, \(\|\mathbf{U}^{n+1}-\mathbf{U}^{n}\|_2/\|\mathbf{U}^{n+1}\|_2\), extrapolated to \(1\ \mu\)s of simulated time, stays below \(10^{-5}\) on several consecutive checks. The steady states are checked a posteriori by the residual of the electron energy equation and by the uniformity of the total current density (Section~\ref{sec:results}).
